% D.84 candidate. UNREGISTERED, NOTES ONLY, pending full independent audits.
\subsection{The local index of a primitive rank-one global translation lattice}

\begin{statement}{D.84}{Open}{Candidate: the whole compact kernel for a primitive rank-one away-five lattice}
Let \(\mathbf u,\mathbf r\) be a basis of \(\mathbb Z_5^2\).
Let \(c=\sum_{q\ne5}c_q\kappa(q)\) be a completed Kummer
class with zero five-valuation whose coefficients generate
\(\mathbb Z_5\). Arbitrary infinite \(c_0\) support is allowed.
Set
\[
 (a,b)=\mathbf u c+\mathbf r\kappa(5),\qquad
 \lambda=\sum_{q\ne5}c_q\frac{\log(q^4)}{4\log6},\qquad
 e=v_5(\lambda),
                                                                  \tag{D.84.1}
\]
where \(v_5(0)=\infty\). For
\(\mathbf r=(r_1,r_2)\), put
\[
 d=\max\{v_5(r_1),v_5(r_2),v_5(r_1-r_2)\}.
                                                                  \tag{D.84.2}
\]
Since \(\mathbf r\) is primitive, either \(d=0\), or exactly
one of these three valuations is positive. Retain the actual
Alexander groups of D.63 and the notation \(Z(n)\) of D.77.
Then
\[
 \mathcal T_s=0,\qquad
 \mathcal S_s\simeq
 \begin{cases}
  \mathbb Z/5,&e>0\text{ and }d=0,\\
  0,&\text{otherwise},
 \end{cases}
 \qquad
 \mathcal K_s\simeq
 \begin{cases}
  0,&e=0,\\
  \mathbb Z/5,&e>0,\ d=0,\\
  Z(\min(e,d)),&e>0,\ d>0.
 \end{cases}                                                     \tag{D.84.3}
\]
The last free factor has countably infinite compact rank when
\(e=d=\infty\). Thus this evaluates a family not covered by
D.82's two independent away-five directions.

The actual point \(t=-24\) has
\[
 (a,b)=(3\kappa(2)+\kappa(3),2\kappa(5)),\qquad
 e=1,\quad d=\infty,\quad \mathcal K_s\simeq\mathbb Z_5.
                                                                  \tag{D.84.4}
\]
Its metabelian fibre is nonempty and is a torsor under
\(\prod_{\mathbb N}\mathbb Z_5\). The theorem makes no
section-existence claim for all the completed parameters,
no general infinitely ramified first-cohomology computation,
and no first higher-derived obstruction computation.
\end{statement}

\begin{proof}
Use \(I_0,k_5,N_1\) from D.82, so
\(N_1=\ker\chi\cap\ker k_5\), the quotient
\(I_0/N_1=\mathbb Z_5(1)\rtimes\Gamma\) is the actual full
affine group, and \(\Gamma=\mathbb Z_5^\times\).
There is an actual element \(k\in I_0\) with
\(\chi(k)=2\) and \(k_5(k)=0\).
Let
\[
 f_{\mathbf w}=t_1^{w_1}t_2^{w_2}-1,\qquad
 P=(t_1-1)(t_2-1)(t_1t_2-1).
\]
The unit statement about the cubic multiplier in D.82 supplies
an actual \(\delta\in N_1\) for which
\(\mathcal B_\delta-1=P E\), \(E\in\Lambda^\times\).
Every element of \(N_1\) has beta multiplier congruent to one
modulo \(P\).

\paragraph{The exact images and the coefficient inclusion.}
D.83's uniformly convergent Kummer cochains and integral local
gauge apply here. The closed away-five lattice is
\(H=\mathbb Z_5\mathbf u\), and the actual two images are
\[
 \begin{split}
 I_s&=\{(h\mathbf u+\mathbf r k_5(g),g):h\in\mathbb Z_5,g\in I_0\},\\
 I_5&=\{(\lambda h\mathbf u+\mathbf r k_5(g),g):
                                      h\in\mathbb Z_5,g\in I_0\}.
 \end{split}                                                     \tag{D.84.5}
\]
Both formulas retain the same actual graph complement. If
\(e=0\), these groups coincide and D.83 proves (D.84.3).
Assume henceforth \(e>0\), allowing infinity.

Put \(D=A_s^D\), and define compact modules and their Cartier
duals by
\[
 B=\Lambda/(f_{\mathbf u}),\qquad
 R=\Lambda/(f_{\lambda\mathbf u}),\qquad
 J=f_{\mathbf u}R,\qquad D_B=B^D,\quad D_R=R^D.
                                                                  \tag{D.84.6}
\]
When \(\lambda=0\), the relation defining \(R\) is absent.
Translation coinduction, with D.81's continuous transversal
when \(\lambda\mathbf u\) is nonprimitive, gives natural
identifications
\[
 H^1(I_s,D)=H^1(I_0,D_B),\qquad
 H^1(I_5,D)=H^1(I_0,D_R).
                                                                  \tag{D.84.7}
\]
If \(\lambda=0\), the second identification simply uses
the graph isomorphism \(I_5\simeq I_0\).
The coefficient action on both compact modules in (D.84.6)
is the actual action
\[
             \rho_{\mathbf r}(g)=
                 t_1^{r_1k_5(g)}t_2^{r_2k_5(g)}\rho_0(g).
                                                                  \tag{D.84.8}
\]
The restriction in (D.84.7) is induced by \(D_B\subset D_R\).
Its entire kernel is therefore
\[
 \mathcal R=
       (J^D)^{I_0}/\operatorname{im}(D_R^{I_0}),
                                                                  \tag{D.84.9}
\]
by the exact sequence \(0\to D_B\to D_R\to J^D\to0\).
The finite affine Chebotarev descent of D.78 gives an injection
\(\mathcal S_s^*\hookrightarrow\mathcal R\). It applies to
each finite globally locally trivial coefficient class before
taking its filtered union, exactly as in D.80 and D.82.

Because \(\mathbf u,\mathbf r\) form a basis, each of the
three boundary linear forms has global coefficient ideal
\(\mathbb Z_5\). Thus D.77 gives
\[
 D_B^{I_0}=D^{I_s}=0,
 \qquad D_R^{I_0}=D^{I_5}\simeq Z(\min(e,d))^*
                       \quad(d>0),                              \tag{D.84.10}
\]
and \(D_R^{I_0}=0\) if \(d=0\).
Here and below a star is ordinary Pontryagin duality after
the Cartier dual has already been taken. In particular
the map from \(D_R^{I_0}\) in (D.84.9) is injective.

\paragraph{The actual compact connecting module.}
On \(N_1\), (D.84.8) is multiplication by beta alone.
The unit cubic element \(\delta\) and the divisibility by
\(P\) for every other element show that
\[
                      (J^D)^{N_1}=(J/PJ)^D.
\]
On this quotient the translation subgroup
\(\mathbb Z_5(1)\subset I_0/N_1\) acts by pure multiplication
\(t_1^{r_1h}t_2^{r_2h}\). Indeed its beta factor is one
on \(J/PJ\). Taking its invariants consequently gives
\[
 (J^D)^{I_0}=(T(-1)_\Gamma)^*,
 \qquad T=J/(PJ+f_{\mathbf r}J).                                 \tag{D.84.11}
\]
The symbol \((-1)\) in this compact formula accounts for
the one Cartier twist: its operator is \(\chi^{-1}\rho\).
Coinvariants are the compact quotient by the closed span of
the operator differences. The \(\Gamma\)-action is the one
induced by actual lifts with \(k_5=0\); differences of lifts
belong to \(N_1\) and have already been killed. Thus (D.84.11)
uses neither a formally assigned beta series nor an assumed
normal coordinate subgroup.

We compute this compact coinvariant module. The actual
permutations of the three punctures, with their transported
frames and the integral gauge in D.79, permit the following
normalization when \(d>0\):
\[
 \mathbf r=(\alpha,\beta),\qquad
 \beta\in\mathbb Z_5^\times,
 \quad v_5(\alpha)=d,
 \quad t_1=t^\beta,\quad t_2=t^{-\alpha},\quad t=1+z.
                                                                  \tag{D.84.12}
\]
The substitution is the exact quotient by \(f_{\mathbf r}\).
The exponent of \(t\) in the image of \(t^{\mathbf u}\)
is \(\beta u_1-\alpha u_2\), a unit. Hence the image of
\(f_{\mathbf u}\) generates \((z)\). Multiplication by its
unit factor merely changes the coordinate of this same
actual ideal. For finite \(n\ge0\), write
\[
 q_n=t^{5^n}-1,\qquad S_n=q_n/z,
 \qquad S_0=1.
                                                                  \tag{D.84.13}
\]
For finite \(e\), the relation obtained after cancelling
the generator \(f_{\mathbf u}\) is \(S_e\), up to a unit.
This remains true for a unit factor in \(\lambda/5^e\):
\((X^\epsilon-1)/(X-1)\) has unit constant coefficient
when \(\epsilon\in\mathbb Z_5^\times\).
For \(e=\infty\) that relation is absent.

In the quotient (D.84.12), \(P\) is a unit times
\(z^2q_d\) when \(0<d<\infty\), it is a unit times
\(z^3\) when \(d=0\), and it is zero when \(d=\infty\).
Thus, as a compact module with its actual action on the
ideal generated by \(z\), (D.84.11) has the scalar description
\[
 T\simeq
 \begin{cases}
  \mathbb Z_5[[z]]/(S_e,z^3),&d=0,\\
  \mathbb Z_5[[z]]/(S_e,z^2q_d),&0<d<\infty,\\
  \mathbb Z_5[[z]]/(S_e),&d=\infty.
 \end{cases}                                                     \tag{D.84.14}
\]
An occurrence of \(S_e\) is omitted for \(e=\infty\).
The scalar description does not discard the generator's
extra linear weight.

\paragraph{The exact finite boundary module.}
For \(n\ge1\), let
\(C_n=\mathbb Z/5^n\), let \(R_n=\mathbb Z_5[C_n]\),
and let \(I_n\) be its augmentation ideal. The multiplication
action \(t\mapsto t^\chi\) gives
\[
                        (I_n^2)_\Gamma\simeq\mathbb Z_5^n.
                                                                  \tag{D.84.15}
\]
Here is an integral verification. The quotient
\(I_n/I_n^2\simeq\mathbb Z/5^n(1)\) has nontrivial
order-four character. Exact order-four coinvariants therefore
identify the coinvariants of \(I_n^2\) and \(I_n\).
The augmentation sequence for \(I_n\subset R_n\) splits
\(\Gamma\)-equivariantly by the group-ring element at zero.
The \(\Gamma\)-orbits in \(C_n\) are zero and the \(n\)
nonzero valuation shells. The coinvariants of the permutation
module \(R_n\) are consequently free on \(n+1\) orbit classes,
and their augmentation kernel is free of rank \(n\).
This proves (D.84.15), with no rational rank inference.

Whenever a quotient in (D.84.14) is
\(Q_n=\mathbb Z_5[[z]]/(S_n)\), its actual Cartier-untwisted
action identifies it with \(I_n^2\). To see the precise map,
regard its scalar element as an element of the ideal generated
by \(f_{\mathbf u}\) and multiply by \(t_1-1\).
Both factors are unit multiples of \(z\). In
\(R_n=\mathbb Z_5[[z]]/(q_n)\), multiplication by their
product identifies \(Q_n\) with \(I_n^2\).
Injectivity follows without dividing by \(5\): if
\(z^2f\in zS_n\mathbb Z_5[[z]]\), then
\(zf=S_nh\); augmentation gives \(5^nh(0)=0\), so
\(h=zh_1\), and cancellation gives \(f=S_nh_1\).
On this quotient \(t_2=1\), so the exact Fox multiplier
is \(\chi S_\chi(t_1)\). Multiplication by \(t_1-1\)
changes it to \(\chi\) times substitution; the Cartier
untwist removes \(\chi\). This proves the asserted
\emph{actual} equivariance with the permutation action.

If \(e\le d\), with finite \(e>0\), then \(q_d\) is
divisible by \(q_e\), or \(P=0\) when \(d=\infty\).
The extra relation in (D.84.14) is redundant. The preceding
comparison and (D.84.15) give
\[
                         T(-1)_\Gamma\simeq\mathbb Z_5^e.
                                                                  \tag{D.84.16}
\]

\paragraph{The three-jet extension when \(0<d<e\).}
Assume \(d\) finite and positive, and \(e>d\), including
\(e=\infty\). There is an exact sequence of compact actual
\(\Gamma\)-modules
\[
                 0\longrightarrow K\longrightarrow T
                       \longrightarrow Q_d\longrightarrow0.
                                                                  \tag{D.84.17}
\]
Before the Cartier untwist, its kernel is
\[
 K\simeq
 \begin{cases}
  (\mathbb Z/5^{e-d})[z]/(z^3),&e<\infty,\\
  \mathbb Z_5[z]/(z^3),&e=\infty,
 \end{cases}                                                     \tag{D.84.18}
\]
with actual generator \(q_d\) in the ideal under discussion.
Indeed \(S_e=S_d g\), where
\(g=\sum_{j=0}^{5^{e-d}-1}t^{j5^d}\), and
\(z^2q_d=z^3S_d\). Cancelling the nonzero element \(S_d\)
in the power-series domain identifies the kernel with
\(\mathbb Z_5[[z]]/(g,z^3)\). Through degree two the
coefficients of \(g\) are all divisible by \(5^{e-d}\),
and its constant quotient is one. This follows directly
by summing \(1,j,\binom j2\), whose denominators two
and six are units at five. Thus \(g\) is
\(5^{e-d}\) times a unit modulo \(z^3\), proving (D.84.18).
The infinite case follows by omitting \(S_e\).

On \(q_d,q_dz,q_dz^2\), the actual compact action has
diagonal characters \(\chi^3,\chi^4,\chi^5\).
The factor \(\chi^2\) is the augmentation of the actual
Alexander multiplier, the factor \(\chi\) comes from
\(q_d(t^\chi)/q_d(t)\), and the remaining powers come
from the displayed monomials. Hence \(K(-1)\) has
diagonal characters \(\chi^2,\chi^3,\chi^4\).
For the actual element with \(\chi=2\), their differences
are units pairwise, so the triangular operator is
diagonalized integrally. The first two eigenvalues minus
one are units; the third, \(2^4-1\), has valuation one.
Since two topologically generates \(\Gamma\),
\[
                         K(-1)_\Gamma\simeq\mathbb F_5,
                                                                  \tag{D.84.19}
\]
generated by the class of \(q_dz^2\).

This possible torsion is killed in the coinvariants of
\(T(-1)\). Set \(g=k^{4\cdot5^{d-1}}\). Its multiplier
is one modulo \(5^d\), so it acts \emph{exactly} as the
identity on \(Q_d(-1)\simeq I_d^2\), by the permutation
description above. Therefore
\((\chi(g)^{-1}\rho_0(g)-1)z^3\) belongs to \(K(-1)\).
Here \(z^3\) is an element of the actual ideal \(J\);
its scalar representative in (D.84.14) is \(z^2\).
Thus the scalar relation \(z^2q_d\) corresponds to
the actual-ideal relation \(z^3q_d\) used below.
D.79's exact actual-power identity gives
\[
 (\chi(g)^{-1}\rho_0(g)-1)z^3
             \equiv 2z^{5^d+2}
                  \pmod{5,\ z^3q_d}.                           \tag{D.84.20}
\]
Changing the primitive coordinate or boundary frame multiplies
the indicated top coefficient by a unit, which suffices here.
Since \(q_dz^2\equiv z^{5^d+2}\pmod5\), this is a nonzero
class of the one-dimensional module (D.84.19). It is an actual
operator difference in \(T(-1)\), so the map from (D.84.19)
to \(T(-1)_\Gamma\) is zero. Right exactness of compact
coinvariants in (D.84.17) now gives
\[
                         T(-1)_\Gamma\simeq\mathbb Z_5^d.
                                                                  \tag{D.84.21}
\]
This step uses an exact identity on the finite boundary
quotient, not merely a congruence of its eigenvalues.

\paragraph{Off the boundary and at the simultaneous infinite limit.}
If \(d=0\), (D.84.14) is just three jets, reduced modulo
\(5^e\) when \(e\) is finite. Its generator has the same
three characters as (D.84.18), so
\[
                         T(-1)_\Gamma\simeq\mathbb F_5.
                                                                  \tag{D.84.22}
\]
There is no finite boundary quotient and no relation of
the kind used in (D.84.20).

If \(e=d=\infty\), the natural map in (D.84.9) must be
checked, rather than comparing two infinite coranks.
Use the actual boundary permutation to take
\(\mathbf r=(0,\beta)\), \(\beta\) a unit.
The translation relation is \(t_2=1\).
The map \(D_R^{I_0}\to(J^D)^{I_0}\), after the natural
axis identifications of D.77 and multiplication by
\(t_1-1\), is dual to the inclusion
\[
 I^2\lhook\joinrel\longrightarrow I,
 \qquad I=(t-1)\mathbb Z_5[[t-1]],
                                                                  \tag{D.84.23}
\]
with pure substitution \(t\mapsto t^\chi\) on both compact
modules. Their quotient \(I/I^2\simeq\mathbb Z_5(1)\)
has nontrivial order-four character. Exact order-four
coinvariants make (D.84.23) an isomorphism on coinvariants,
and hence on full \(\Gamma\)-coinvariants. Its invariant
dual map is consequently an isomorphism. This proves
\(\mathcal R=0\) also in this infinite case.

For finite \(\min(e,d)>0\), (D.84.16)--(D.84.21) show
that \((J^D)^{I_0}\) is a divisible group of that finite
corank. The injective map in (D.84.10) has source a
divisible group of the same corank, so it is an isomorphism.
Indeed its Pontryagin dual is a surjection between free
\(\mathbb Z_5\)-modules of the same finite rank, hence an
isomorphism. Together with (D.84.22) and
\(D_R^{I_0}=0\) when \(d=0\), this computes the whole
restriction kernel:
\[
 \mathcal R\simeq\mathbb F_5\quad(e>0,d=0),
 \qquad \mathcal R=0\quad\text{otherwise}.
                                                                  \tag{D.84.24}
\]

\paragraph{The nonzero class and a common finite global level.}
Let \(c_3\) denote D.78's normalized raw cubic cocycle
modulo five, viewed in the constant Cartier submodule
\(\mathbb F_5(-1)\subset D_B\).
Its global class is nonzero in the present, nonconstant
coinduced coefficient. To prove this directly, choose
an actual \(\tau\in I_0\) with
\(\chi(\tau)=1,k_5(\tau)=1,c_3(\tau)=0\).
Such an element exists: first prescribe the two affine
coordinates, then multiply by a suitable element of
\(N_1\), on which the cubic character is primitive.
On \(B\simeq\mathbb Z_5[[z]]\), its compact operator is
multiplication by \(t^{\mathbf r}\mathcal B_\tau\).
The difference from one is \(z\) times a unit, because
\(\mathbf u,\mathbf r\) are a basis and beta starts in
degree three on the cyclotomic kernel.
Thus the \(\tau\)-fixed part of \(D_B\) is precisely the
constant functional module. If \(c_3\) were the coboundary
of an element of \(D_B\), that element would be fixed
by \(\tau\), hence constant. An element of \(N_1\) with
unit cubic character fixes constants and contradicts
the coboundary identity. This proves nonvanishing in
\(H^1(I_0,D_B)\), hence by coinduction and inflation in
\(H^1(G_{\mathbb Q},D)\).

Suppose \(e>0,d=0\). Modulo five the local unit vector
\(\lambda\mathbf u\) is zero, and the valuation vector
\(\mathbf r\) is a nonboundary direction. D.82's actual
four-dimensional line quotient of the fixed cubic module
therefore gives a local primitive for \(c_3\) at five.
This local construction uses only the residual vectors
and the actual beta coefficients; it does not require
D.82's global away-five rank-two hypothesis.

At every prime away from five with nonzero abelian
inertia, D.80's finite local Tate-duality argument and
D.43 give \(H^1(G_{\mathbb Q_q},D)=0\).
Only finitely many such primes have nonzero residual
inertia modulo five. At those primes a primitive occurs
in some finite coefficient module, and one takes a
common larger level for this finite set. At every
other prime, the original ten-dimensional cubic module
is unramified and D.78's explicit Frobenius primitive
works at the same fixed level. The real place contributes
no odd-primary cohomology. This is precisely D.80's
finite-exception argument; it gives a single finite
global coefficient module in which \(c_3\) is locally
trivial everywhere. Its nonzero image just proved gives
\(\mathcal S_s^*\ne0\). The upper bound (D.84.24)
then proves the finite assertion in (D.84.3).

\paragraph{The whole compact kernel and actual examples.}
The global three boundary depths are zero. If \(d=0\),
the local valuation vector alone makes every local
depth zero. If \(d>0\), there is one boundary form
whose value on \(\mathbf r\) has valuation \(d\);
its value on \(\mathbf u\) is a unit because these
vectors are a basis. Its local depth is therefore
\(\min(e,d)\), and the other two are zero. D.77 gives
all the remaining conclusions of (D.84.3), continuously.

For \(t=-24\), the global Kummer pair is exactly
(D.84.4), since signs vanish in pro-five completion.
Its away-five coefficients \(3,1\) generate \(\mathbb Z_5\),
and the basis vectors are \((1,0),(0,2)\).
The exact identity
\[
                      24^4-1=25\cdot13271
\]
shows that its unit logarithm coefficient has valuation
one. The actual full point section supplies existence;
D.74 then computes its nonempty metabelian fibre.

For explicit completed modules at every index, let
\(\lambda_2,\lambda_3\) be the units in D.78 and set
\[
 c_m=\kappa(2)+\frac{5^m-\lambda_2}{\lambda_3}\kappa(3)
       \quad(m\ge1),\qquad
 c_\infty=\kappa(2)-\frac{\lambda_2}{\lambda_3}\kappa(3).
                                                                  \tag{D.84.25}
\]
These classes are globally primitive and have the stated
local coefficient \(5^m\), or zero.
Choosing \(\mathbf u=(1,0),\mathbf r=(0,1)\) realizes
\(Z(m)\), including the infinite product, while
\(\mathbf r=(1,2)\) realizes \(\mathbb Z/5\).
They are actual finite-support Galois modules; no further
section existence is inferred.
\end{proof}
